% Cote 119: offprints, reports, courses, CV and letters (typed and printed part): archive cote 119, site page 297 (PDF page 298).
% typescript. Draft compared: none.
% AI-transcribed by Claude Opus 5.5 (claude-opus-5-5, adaptive thinking, effort high), Message Batches API, batch r119_full (msgbatch_01DwZNm5Y7YsKBWyAN4sFs8w), 2026-11-13; not yet checked by eye.
\scanpage{119/297}
\mspage{- 10 -}

De ce critère, combiné avec (4.4), Bayer-Lenstra déduisent le théorème
suivant :

(6.1) - \uline{S'il existe une extension de degré impair de} $k$ \uline{sur laquelle} $L$
\uline{acquiert une} BNA, \uline{alors} $L$ \uline{a une} BNA \uline{sur} $k$.
En particulier :
(6.2) - \uline{Si} $G$ \uline{est d'ordre impair}, \uline{toute} $G$-\uline{algèbre galoisienne a une} BNA.

Voici quelques autres résultats relatifs aux BNA ; les démonstrations
seront publiées en collaboration avec E.Bayer.

Soit $L$ une $G$-algèbre galoisienne, et soit $\hand{\varphi}_L : \mathrm{Gal}(k_s/k) \;{-}{>}\; G$
l'homomorphisme correspondant. Si $x$ est un élément de $H^n(G, Z/2Z)$,
son image par $\hand{\varphi}_L^{*} : H^n(G, Z/2Z) \;{-}{>}\; H^n(\mathrm{Gal}(k_s/k), Z/2Z) = H^n(k, Z/2Z)$
sera notée $x_L$ .

(6.3) \uline{Pour que} $L$ \uline{ait une} BNA, \uline{il faut que} $x_L = 0$ \uline{pour tout} $x \hand{\in} H^1(G, Z/2Z)$
(autrement dit, l'image de $\mathrm{Gal}(k_s/k)$ dans $G$ doit être contenue dans tous
les sous-groupes d'indice $2$ de $G$). \uline{Cette condition est suffisante si la}
2-\uline{dimension cohomologique de} $\mathrm{Gal}(k_s/k)$ \uline{est} $\hand{\leqslant}$ 1 (autrement dit si les
2-sous-groupes de Sylow de $\mathrm{Gal}(k_s/k)$ sont des pro-2-groupes libres).

(6.4) - \uline{Supposons que} $k$ \uline{soit un corps de nombres}. \uline{Pour que} $L$ \uline{ait une}
BNA, \uline{il faut que} $\hand{\varphi}_L(c_v) = 1$ \uline{pour toute place réelle} $v$ \uline{de} $k$ ($c_v$ dési-%
gnant la conjugaison complexe relative à une extension de $v$ à $k_s$). \uline{Cette}
\uline{condition est suffisante si} $H^1(G, Z/2Z) = H^2(G, Z/2Z) = 0$.

(6.5) \uline{Le cas où un} 2-\uline{groupe de Sylow de} $G$ \uline{est abélien élémentaire}
Soit $S$ un 2-sous-groupe de Sylow de $G$. Supposons que $S$ soit un groupe
abélien élémentaire d'ordre $2^r$ , $r \hand{\geqslant} 1$; l'ordre de $G$ est $2^r m$ , avec $m$
impair.
(6.5.1)-\uline{Il existe une} $r$-\uline{forme de Pfister} $q^1_L$ , \uline{et une seule à isomorphisme}
\uline{près}, \uline{telle que} $2^r q_L \hand{\simeq} m \hand{\otimes} q^1_L$ (somme directe de $m$ copies de $q^1_L$).

Cette forme constitue un \uline{invariant} de l'algèbre galoisienne $L$ considérée.
C'est la forme unité si $L$ a une BNA. Réciproquement :

(6.5.2)- \uline{Supposons que le normalisateur} $N$ \uline{de} $S$ \uline{opère transitivement sur}
$S - \hand{\{}1\hand{\}}$ . \uline{Il y alors équivalence entre} :
(i) $L$ \uline{a une} BNA.
(ii) \uline{La forme} $q_L$ \uline{est isomorphe à la forme unité de rang} $2^r m$ .
(iii) \uline{La forme} $q^1_L$ \uline{est isomorphe à la forme unité de rang} $2^r$.

\margin{297}
